\chapter*{Hills and Mountains}\stepcounter{chapter}\phantomsection\addcontentsline{toc}{chapter}{Hills and Mountains}

\section*{Duergar's Cottage}\phantomsection\addcontentsline{toc}{section}{Duergar's Cottage}
{\entryfont Psychological Hazard (Illusion-Based, Optional Combat)\\
\textbf{Recommended Level:} 5 - 8\\
\textbf{Difficulty:} Moderate (High if players act on false assumptions)}
\subsection*{Setup}
{\entryfont As the party traverses the hills or mountainous terrain - likely at dusk or night - they may already feel disoriented, with the landscape repeating itself in unsettling ways. At an opportune moment, they spot a warm, inviting light in the distance: a small cottage that appears impossibly well-placed in such a hostile environment.}
\begin{DndReadAloud}
	The hills seem to stretch endlessly around you, their shapes softened and distorted by creeping mist. Landmarks that once felt certain now blur together, and the path - if there ever was one - has all but vanished beneath your feet.

	The air grows colder as the light fades, and an uneasy sense settles in: you may no longer be entirely sure where you are.

	Then, through the haze, a flicker - warm, steady, unmistakable.

	A light.

	Drawing closer, the outline of a small cottage emerges, its windows glowing softly against the gloom. Smoke curls lazily from a chimney, and the promise of warmth pushes back the encroaching chill.

	Inside, a fire crackles invitingly. Thick wooden beams lie stacked nearby, and the space feels sheltered, quiet... safe.
\end{DndReadAloud}

\subsection*{Mechanics}
{\entryfont This is not a conventional illusion. The cottage is created via Mirage Arcane and is fully tangible - objects behave exactly as expected, and no checks reveal the illusion. Only Truesight exposes that something is fundamentally wrong.

The danger lies in a specific inconsistency: the floor directly behind the fire, where the thick wooden beams lie, is not real. This section is concealed by Hallucinatory Terrain, masking open air over an 80 ft drop.

\begin{itemize}
	\item A DC 14 Intelligence (Investigation) check reveals the inconsistency
	\item A creature stepping onto this area falls immediately (\DndDice{8d6} Bludgeoning damage)
\end{itemize}

\noindent The Duergar enters only after the party settles. It remains silent, breaks a beam to feed the fire, and fixes a player with an unblinking stare - clearly inviting imitation without ever stating it.}

\subsection*{Possible Solutions}
{\entryfont Encourage players to approach the situation with caution, curiosity, or restraint rather than assumption.}
\subsubsection*{Physical caution}
{\entryfont Players who probe the environment - testing the ground with tools, throwing objects, or avoiding blind movement - can reveal the dangerous section without triggering it. Maintaining distance from the far side of the fire is often sufficient.}
\subsubsection*{Magical Awareness}
{\entryfont While Mirage Arcane itself does not give itself away, Detect Magic reveals a strong illusion aura permeating the entire area, hinting that something is fundamentally wrong even if not precisely where. Dispel Magic can suppress or remove parts of the illusion (at your discretion, potentially collapsing the false floor or destabilizing the environment), immediately exposing the hazard. Truesight cleanly reveals both layers of deception.}
\subsubsection*{Behavioural restraint}
{\entryfont If the players refuse to engage with the Duergar's silent challenge - neither mimicking its actions nor reacting to its stare - the creature eventually loses interest. After a period of inactivity, it may simply leave, ending the encounter without combat or escalation.}

\subsection*{Possible Complications}
{\entryfont The encounter escalates due to the player's psychology and paranoia rather than outside influences.
\begin{itemize}
	\item If the party agitates the \hyperref[monster:Duergar]{\LinkFont{Duergar}} - whether by openly challenging it, disrupting its "game", or casting Dispel Magic in its presence - the creature immediately shifts from passive observer to active threat, responding with aggression and likely initiating combat.
	\item At the same time, the nature of the illusion can create tension within the party itself. Differing interpretations of what is real or safe may lead to hesitation, conflicting decisions, or impulsive actions, increasing the likelihood that someone takes a dangerous risk.
\end{itemize}}

\subsection*{Outcomes}
\subsubsection*{Endurance}
{\entryfont If the party refuses to engage and avoids the hazard, the Duergar eventually loses interest. After some time - or by dawn - it withdraws, and the illusion fades, revealing the cliff. The party gains no material reward but succeeds through caution and restraint.}
\subsubsection*{Failure}
{\entryfont If a character steps onto the false floor, they fall 80 feet as the illusion gives way. The Duergar silently observes before disappearing, leaving the party to deal with the aftermath as the illusion soon collapses.}
\subsubsection*{Duergar Defeated}
{\entryfont If the Duergar is slain, the illusion lingers briefly before fading. The cliff becomes fully visible, and the party can safely assess the area and claim any loot it carried.
\begin{itemize}
	\item Pouch containing 80 GP
	\item Magical Mote to cast Major Image once
\end{itemize}}

\clearpage

\section*{Spriggan's Barrow}\phantomsection\addcontentsline{toc}{section}{Spriggan's Barrow}
{\entryfont Burial Mound with Territorial Fey Guardians\\
\textbf{Recommended Level:} 5 - 7\\
\textbf{Difficulty:} Hard}
\subsection*{Setup}
{\entryfont A windswept burial mound rises from the rolling hills of Caledonia, marked by ancient stones half-sunken into the earth. Travellers may stumble upon it while crossing the highlands or deliberately seek it out due to rumours of buried relics. The barrow is sacred to a trio of Spriggans who silently observe intruders from hidden vantage points. They do not reveal themselves unless the mound is disturbed - digging, tampering with stones, or removing artifacts is seen as a grave violation, prompting immediate retaliation.}
\begin{DndReadAloud}
	The wind quiets as you approach the hilltop, where a great mound of earth rises from the land, encircled by leaning stones worn smooth by age. The grass here grows thick and uneven, and the air feels heavy, as though weighed down by something unseen. Faint impressions mark the soil - old disturbances, long since settled - and a lingering stillness presses in around you, broken only by the distant cry of wind across the hills.
\end{DndReadAloud}

\subsection*{Mechanics}
{\entryfont The encounter consists of three \hyperref[monster:Spriggan]{\LinkFont{Spriggans}}. They begin in Small Form, emerging only after the barrow is disturbed. One may initially remain hidden to flank or ambush. Early in combat, at least one Spriggan uses Spriggan's Growth to become Large and engage in melee, while the others alternate between ranged Boulder attacks and evasive tactics using Nimble Escape. They prioritize targets interacting with the burial site and fight until only one remains, which may attempt to retreat.}

\subsection*{Possible Solutions}
\subsubsection*{Observation \& Restraint}
{\entryfont Careful examination of the mound (Perception / Insight) reveals subtle signs of guardianship - disturbed grass, unnatural silence, or fleeting movement - suggesting the site is protected. Choosing not to disturb the barrow avoids conflict entirely.}
\subsubsection*{Magical Awareness}
{\entryfont Spells such as Detect Magic reveal faint transmutation and enchantment traces tied to the Spriggans' presence. Dispel Magic can interfere with their environmental manipulation but does not remove their territorial aggression once triggered.}
\subsubsection*{Respectful Conduct}
{\entryfont If the party treats the site with reverence - leaving offerings, avoiding excavation, or verbally acknowledging the dead - the Spriggans may choose not to engage, remaining unseen and allowing passage.}
\vfill\eject
\subsection*{Possible Complications}
\subsubsection*{Ambush \& Isolation}
{\entryfont The Spriggans exploit terrain and timing, striking when the party is divided or distracted. Fear effects and hit-and-run tactics isolate weaker targets, forcing the party into unfavourable positioning.}
\subsubsection*{Escalation Through Interference}
{\entryfont Attempting to dispel magic or aggressively search the mound in the Spriggans' presence immediately escalates the encounter into full hostility, removing any chance of avoidance.}
\subsubsection*{Party Conflict}
{\entryfont Disagreements within the group - whether to respect the dead or loot the barrow - can delay decisions and leave individuals exposed, giving the Spriggans time to prepare a decisive strike.}

\subsection*{Outcomes}
{\entryfont If the party leaves the barrow undisturbed, nothing occurs beyond an uneasy stillness - the Spriggans remain hidden, silently observing but choosing not to act. However, if the barrow is disturbed and conflict with the Spriggans ensues, the outcome shifts and is instead determined by the result of that confrontation.}
\subsubsection*{Spriggans Retreat}
{\entryfont If at least one Spriggan escapes, the burial mound begins to shift and settle unnaturally, as if shaped by ancient magic akin to Move Earth. The soil folds inward, stones grind and sink, and any exposed burial contents are swallowed and hidden deep beneath the mound, becoming effectively unrecoverable. Only a few overlooked remnants remain scattered near the surface:
\begin{itemize}
	\item A single carved idol (25 GP)
	\item 1d4 × 5 GP in old, weathered coinage
	\item A simple but well-crafted trinket (10 GP)
\end{itemize}}
\subsubsection*{Spriggans are Defeated}
{\entryfont With no guardians remaining, the barrow can be excavated without interference. The party uncovers a modest cache of ancient grave goods:
\begin{itemize}
	\item 1d4 carved stone idols (25 gp each)
	\item A bronze jewellery piece (75 gp)
	\item \hyperref[mi:OvergrownBarkshield]{\LinkFont{Overgrown Barkshield}}
\end{itemize}}
\subsubsection*{Party Retreats}
{\entryfont The Spriggans reclaim the mound and remain vigilant. Any items left behind are buried or removed, and the site becomes overtly hostile to the party upon their return.}